\subsection{扩张的代数自由族}

定义6. --- 设 L 是域 K 的一个扩域， \((\mathrm{E}_{i})_{i\in I}\) 是 L 的子扩张的一个族。族 \((\mathrm{E}_{i})_{i\in I}\) 称为代数自由的，如果下列条件满足：

(AF) 对每个 \(i \in I\) ，设 \(\mathbf{A}_i\) 是 \(\mathbf{E}_i\) 的一个在 \(\mathbf{K}\) 上代数自由的子集。则 \(\mathbf{A}_i \cap \mathbf{A}_j = \emptyset\) （对 \(i \neq j\) ），且 \(\cup \mathbf{A}_i\) 在 \(\mathbf{K}\) 上代数自由。

注记. --- 由命题3（V，第107页），只需对有限子集 \(A_{i}\) 验证条件（AF）即可。由此我们得到如下结果：若 \((\mathrm{E}_{i})_{i\in I}\) 是代数自由族，则 \((\mathrm{E}_{i}^{\prime})_{i\in I}\) 亦然，只要 \(E_{i}^{\prime}\) 是 \(E_{i}\) 的一个子扩张（对每个 \(i\in I\) ）；反之，若每个族 \((\mathrm{E}_{i}^{\prime})_{i\in I}\) （其中 \(E_{i}^{\prime}\) 是 E 的一个有限生成子扩张，对每个 \(i\in I\) ）都是代数自由的，则 \((\mathrm{E}_{i})_{i\in I}\) 是代数自由的。另一方面，为使 \((\mathrm{E}_{i})_{i\in I}\) 是代数自由的，充要条件是 \((\mathrm{E}_{i})_{i\in J}\) 对 I 的每个有限子集 J 都是代数自由的。直观地可以说，扩张的代数无关性是一种“有限特征”的性质。

命题15. --- 设 \((\mathbf{E}_i)_{i\in I}\) 是域 K 的一个给定扩张 L 的子扩张族。则下列条件等价：

\begin{enumerate}
\def\labelenumi{\alph{enumi})}
\item
  族 \((\mathrm{E}_i)_{i\in I}\) 是代数自由的。
\item
  对每个 \(i \in I\) ，扩张 \(E_i\) 在 \(K\) 上与扩张 \(F_i\) （由诸 \(E_j\) （ \(j \neq i\) ）生成的）代数不相交。
\item
  存在一个族 \((\mathbf{B}_i)_{i\in I}\) （由 \(L\) 的两两不相交的子集组成），使得 \(\mathbf{B}_i\) 是 \(\mathbf{E}_i\) 在 \(K\) 上的一个超越基（对每个 \(i\in I\) ），且 \(\mathbf{B} = \bigcup_{i\in I}\mathbf{B}_i\) 在 \(K\) 上代数自由。
\end{enumerate}

显然 \(a)\) 蕴含 \(c)\) 。

假设 \(c\) )，我们选取 \(i\) （属于 \(I\) ）；令 \(C_i = \bigcup_{j \neq i} B_j\) 。对每个 \(j \neq i\) ，每个元素

（属于 \(E_{j}\) 的）在 \(\mathbf{K}(\mathbf{B}_{j})\) 上是代数的，从而在 \(\mathbf{K}(\mathbf{C}_{i})\) 上更是如此。由 V 第18页的推论1，域 \(F_{i}\) 因此在 \(\mathbf{K}(\mathbf{C}_{i})\) 上是代数的。此外，我们有 \(B_{i} \cap C_{i} = \varnothing\) 且 \(B = B_{i} \cup C_{i}\) 在 K 上代数自由；因此 \(B_{i}\) 在 \(\mathbf{K}(\mathbf{C}_{i})\) 上代数自由（V，第107页，命题4），从而也在 \(F_{i}\) 上代数自由（它在 \(\mathbf{K}(\mathbf{C}_{i})\) 上是代数的），由 V 第108页的命题6。于是我们证明了 \(E_{i}\) 在 K 上与 \(F_{i}\) 代数不相交（V，第113页，命题12），故 c) 蕴含 b)。

现在假设 \(b)\) 并证明 \(a)\) 。只需证明：若 \(i_1, \ldots, i_n\) 是 \(I\) 中互不相同的元素，则扩张族 \((E_{i_1}, \ldots, E_{i_n})\) 是代数自由的；我们对 \(n\) 作归纳论证，情形 \(n = 1\) 是平凡的。现设 \(n > 1\) 且族 \((E_{i_1}, \ldots, E_{i_{n-1}})\) 是代数自由的；对 \(1 \leqslant k \leqslant n\) ，选取一个子集 \(A_k\) （ \(E_{i_k}\) 的、在 \(K\) 上代数自由的），并设

\(B = A_{1} \cup \ldots \cup A_{n-1}\) 。由归纳假设，子集 \(A_{1}, \ldots, A_{n-1}\) 两两不相交，且 B 在 K 上代数自由；由 b)， \(E_{i_{n}}\) 与 \(F_{i_{n}}\) 代数不相交，且由于 B 包含于 \(F_{i_{n}}\) ，我们有 \(B \cap A_{n} = \varnothing\) 且 \(B \cup A_{n} = A_{1} \cup \ldots \cup A_{n}\) 在 K 上代数自由。于是我们证明了族 \((E_{i_{1}}, \ldots, E_{i_{n}})\) 是代数自由的。

下面的命题推广了命题13（V，第113页）的 \(b)\) 部分。

命题16. --- 设 \((\mathrm{E}_i)_{i\in I}\) 是域 K 的一个扩张的子扩张族，并设 E 是由 \(\cup\) \(\mathrm{E}_i\) 生成的域。

\begin{enumerate}
\def\labelenumi{\alph{enumi})}
\item
  我们有 \(\operatorname{tr} \cdot \deg_{\mathbf{K}} \mathrm{E} \leqslant \sum_{i \in I} \operatorname{tr} \cdot \deg_{\mathbf{K}} \mathrm{E}_i\) ，并且当族 \((\mathrm{E}_i)_{i \in I}\) 在 \(\mathbf{K}\) 上代数自由时取等号。
\item
  反之，假设 \(\operatorname{tr} \cdot \deg_{\mathbf{K}} \mathrm{E} = \sum_{i \in I} \operatorname{tr} \cdot \deg_{\mathbf{K}} \mathrm{E}_i\) 且 \(\operatorname{tr} \cdot \deg_{\mathbf{K}} \mathrm{E}\) 有限；则族 \((\mathrm{E}_i)_{i \in I}\) 在 \(\mathbf{K}\) 上代数自由。
\end{enumerate}

对每个 \(i \in I\) ，设 \(\mathbf{B}_i\) 是 \(\mathbf{E}_i\) 在 \(\mathbf{K}\) 上的一个超越基，并设 \(\mathbf{B} = \bigcup_{i \in I} \mathbf{B}_i\) 。对每个 \(i \in I\) ， \(\mathbf{E}_i\) 的每个元素在 \(\mathbf{K}(\mathbf{B}_i)\) 上是代数的，从而在

K(B) 上是代数的。现在 V 第18页的推论1表明 E 在 K(B) 上是代数的；由 V 第109页定理2，B 包含 E 在 K 上的一个超越基。此外，若族 \((\mathrm{E}_{i})_{i\in I}\) 在 K 上是代数自由的，则诸 \(B_{i}\) 两两不相交，且集合 B 在 K 上代数自由。这就建立了 a)（集合论，III，第160页，命题4的推论）。

在 \(b\) ) 的假设下，E 在 \(\mathbf{K}(\mathbf{B})\) 上是代数的，且在 \(\mathbf{K}\) 上超越次数有限，并且我们有 \(\operatorname{Card}(\mathbf{B}) \leqslant \operatorname{tr} \cdot \deg_{\mathbf{K}} \operatorname{E}\) 。由 V 第110页的推论1，B 在 \(\mathbf{K}\) 上代数自由，且诸 \(\mathbf{B}_i\) 两两不相交。现在命题15表明族 \((\mathrm{E}_i)_{i \in I}\) 在 \(\mathbf{K}\) 上代数自由。

在陈述下一个定理之前，我们指出存在具有任意超越次数的 K 的代数闭扩张，例如一个适当的有理分式域的代数闭包。

定理5. --- 设 \((\mathrm{E}_i)_{i\in I}\) 是域 \(\mathbf{K}\) 的一族扩张， \(\Omega\) 是 \(\mathbf{K}\) 的一个代数闭扩张。假设不等式

\[
\operatorname{tr} \cdot \deg_ {K} \Omega \geqslant \sum_ {i \in I} \operatorname{tr} \cdot \deg_ {K} E _ {i}\tag{2}
\]

成立。则存在一个代数自由族 \((\mathrm{F}_i)_{i\in I}\) （由 \(\Omega\) 的子扩张组成），使得 \(\mathbf{F}_i\) 与 \(\mathbf{E}_i\) 是 K-同构的（对所有 \(i\in I\) ）。

对每个 \(i \in I\) ，设 \(B_i\) 是 \(E_i\) 在 \(K\) 上的一个超越基。设 \(B\) 是 \(\Omega\) 在 \(K\) 上的一个超越基。由 (2) 我们有 Card \(B \geqslant \sum_{i \in I} \text{Card } B_i\) ；因此

存在 B 的两两不相交子集的一个族 \((\mathbf{B}_{i}^{\prime})_{i\in I}\) 以及双射 \(u_{i}:B_{i}\to B_{i}^{\prime}\) （对 \(i\in I\) ）。由 V 第106页的命题2， \(u_{i}\) 扩张为一个 K-同构 \(v_{i}\) （从 \(\mathbf{K}(\mathbf{B}_{i})\) 到 \(\mathbf{K}(\mathbf{B}_{i}^{\prime})\) 上的）；由于 \(\Omega\) 是代数闭的且 \(E_{i}\) 在

\(K(B_{i})\) 上是代数的，推论（V，第23页）表明 \(v_{i}\) 延拓为 \(E_{i}\) 到一个子扩张 \(F_{i}\) （ \(\Omega\) 的）上的一个 K-同构。根据构造， \(B_{i}^{\prime}\) 是 \(F_{i}\) 在 K 上的一个超越基，且命题15（V，第115页）表明子扩张族 \((F_{i})_{i\in I}\) （ \(\Omega\) 的）在 K 上是代数自由的。

推论1. --- 设 E 和 Ω 是域 K 的两个扩张。假设 Ω 是代数闭的，其超越次数至少等于 E 的超越次数。则 E 与 Ω 的一个子扩张是 K-同构的。

推论2. --- 设 \(\Omega\) 是一个在其素子域上具有无限超越次数的代数闭域。则每个与 \(\Omega\) 特征相同的域都是同构于 \(\Omega\) 的子域的域的上升有向并。

因为每个域都是其素域上有限生成的子域的上升有向并，于是只需应用推论1即可。

* 例. --- 这特别适用于特征 0 的情形，即取 \(\Omega = \mathbf{C}\) （“莱夫谢茨原理”）。 *

